% !TEX TS-program = pdflatex
% !TEX encoding = UTF-8 Unicode

% This file is a template using the "beamer" package to create slides for a talk or presentation
% - Giving a talk on some subject.
% - The talk is between 15min and 45min long.
% - Style is ornate.

% MODIFIED by Jonathan Kew, 2008-07-06
% The header comments and encoding in this file were modified for inclusion with TeXworks.
% The content is otherwise unchanged from the original distributed with the beamer package.

\documentclass{beamer}


% Copyright 2004 by Till Tantau <tantau@users.sourceforge.net>.
%
% In principle, this file can be redistributed and/or modified under
% the terms of the GNU Public License, version 2.
%
% However, this file is supposed to be a template to be modified
% for your own needs. For this reason, if you use this file as a
% template and not specifically distribute it as part of a another
% package/program, I grant the extra permission to freely copy and
% modify this file as you see fit and even to delete this copyright
% notice. 


\mode<presentation>
{
  \usetheme{Warsaw}
  % or ...

  \setbeamercovered{transparent}
  % or whatever (possibly just delete it)
}


\usepackage[english]{babel}
% or whatever

\usepackage[utf8]{inputenc}
% or whatever

\usepackage{mathrsfs}

\usepackage{times}
\usepackage[T1]{fontenc}
% Or whatever. Note that the encoding and the font should match. If T1
% does not look nice, try deleting the line with the fontenc.

%%% MATH RELATED
\usepackage{amsmath,amsthm,amssymb}
\newtheorem{prob}{Problem}
\newtheorem{idea}{Idea}
\newtheorem{defn}{Definition}
\newtheorem{thm}{Theorem}
\newtheorem{conj}{Conjecture}
\newtheorem{ex}{Example}

\newcommand{\bbc}{\mathbb{C}}
\newcommand{\bbr}{\mathbb{R}}
\newcommand{\bbn}{\mathbb{N}}
\newcommand{\Ai}{\text{Ai}}
\newcommand{\Be}{\text{Be}}
\newcommand{\wt}[1]{\widetilde{#1}}

\title{Math GRE Prep Day 4} 
\subtitle
{} % (optional)

\author[W.R. Casper] % (optional, use only with lots of authors)
{W.R. Casper}
% - Use the \inst{?} command only if the authors have different
%   affiliation.

\institute[California State University Fullerton] % (optional, but mostly needed)
{
  Department of Mathematics\\
  California State University Fullerton}
% - Use the \inst command only if there are several affiliations.
% - Keep it simple, no one is interested in your street address.

\subject{Talks}
% This is only inserted into the PDF information catalog. Can be left
% out. 



% If you have a file called "university-logo-filename.xxx", where xxx
% is a graphic format that can be processed by latex or pdflatex,
% resp., then you can add a logo as follows:

% \pgfdeclareimage[height=0.5cm]{university-logo}{university-logo-filename}
% \logo{\pgfuseimage{university-logo}}



% Delete this, if you do not want the table of contents to pop up at
% the beginning of each subsection:
\AtBeginSubsection[]
{
  \begin{frame}<beamer>{Outline}
    \tableofcontents[currentsection,currentsubsection]
  \end{frame}
}


% If you wish to uncover everything in a step-wise fashion, uncomment
% the following command: 

%\beamerdefaultoverlayspecification{<+->}


\begin{document}

\begin{frame}
  \titlepage
\end{frame}

\begin{frame}{Outline}
  \tableofcontents
  % You might wish to add the option [pausesections]
\end{frame}

% Since this a solution template for a generic talk, very little can
% be said about how it should be structured. However, the talk length
% of between 15min and 45min and the theme suggest that you stick to
% the following rules:  

% - Exactly two or three sections (other than the summary).
% - At *most* three subsections per section.
% - Talk about 30s to 2min per frame. So there should be between about
%   15 and 30 frames, all told.

\section{Calculus III Problems}

\subsection{Basic problems}

\begin{frame}{Problem 1}
Find the Jacobian of the transformation 
$$\left\lbrace\begin{array}{cc}
u &= xe^{xy}\\
v &= ye^{xy}\\
\end{array}\right.$$
\end{frame}

\begin{frame}{Problem 2}
Let $f(x,y) = x^3-axy+y^2-x$.
For what values of $a$ does the function have a local minimum?
\end{frame}

\begin{frame}{Problem 3}
Find $$\iint_R z dxdy$$ where $z=8xy$ and $R$ is the part of the unit disk in the first quadrant.
\end{frame}

\begin{frame}{Problem 4}
Find the length of the parametric curve $x = \cos(t)e^{t}$ and $y = -\sin(t)e^t$ between $0<t<1$.
\end{frame}

\begin{frame}{Problem 5}
If $\mathbf F = \langle 2y-2x,x^2+y\rangle$, find the value of $\int_C\mathbf{F}\cdot d\mathbf{r}$ where $C$ is the segment of the parabola $y=x^2$ starting from $(-1,1)$ and going to $(0,0)$.
\end{frame}

\begin{frame}{Problem 6}
Let $C$ be the portion of the astroid $x^{2/3}+y^{3/2}=1$ from $(1,0)$ to $(0,1)$, which is parameterized by the equations
$$x=\cos^3 t,\quad y=\sin^3 t,\ \ 0\leq t\leq \pi/2.$$
Evaluate the integral
$$\int_C(y\cos(xy)-1)dx + (1+x\cos(xy))dy$$
\end{frame}

\begin{frame}{Problem 7}
Find all functions $f(x,y)$ satisfying $\frac{\partial f}{\partial x} = 2x + y$ and $\frac{\partial f}{\partial x} = x + 2y$.
\end{frame}



\begin{frame}{Problem 8}
Find the point on the plane $2x + y + 3z = 3$ which is closest to the origin.
\end{frame}

\begin{frame}{Problem 9}
Set up an integral which represents the volume of the solid bounded above
by the graph of $z = 6 - x^2 - 2y^2$ and below by the graph of $z = -2 + x^2 + 2y^2$.
\end{frame}

\begin{frame}{Problem 10}
Minimize the function $f(x, y, z) = x + 4z$ on the surface $x^2+y^2+z^2=1$.
\end{frame}

\end{document}


